\section{基础设施谱系：HPC、Cloud、Data Cloud 到 AI Cloud}

\subsection{从科学计算到通用云}
讲者回顾了早期 HPC 集群时代：高算力任务（物理、天气等）运行在集中式集群上，强调并行计算能力。随后 VPS 与公有云出现，把硬件运维抽象出去，让开发者以更低门槛部署业务。

\begin{knowledgebox}{公有云成功的抽象前提}
机器被视作可替换资源，工作负载多为松耦合服务。只要接口一致，底层物理位置通常不影响业务逻辑。
\end{knowledgebox}

\subsection{Web 2.0 到 Data Cloud}
通用云解决了机器供给和服务托管后，下一轮瓶颈转向海量行为数据的存储、处理与分析。本小节承接官方基础设施时间轴，说明 Snowflake、Databricks 一类 data cloud 为什么从 web workload 中长出，并为后面的 AI 数据面问题提供历史参照。
随着用户行为数据激增，系统重心从 ``托管服务'' 走向 ``数据分析与推荐''。讲者以 Snowflake 与 Databricks 为代表，指出 Data Cloud 时代核心是大规模数据处理与 SQL/批流计算统一。

\subsection{AI Cloud 的新矛盾：IO 与数值计算}
Data Cloud 主要优化数据移动与查询，AI Cloud 则让大规模数值计算、模型状态和跨卡通信同时成为核心。本小节把 I/O 与 compute 放进同一资源模型，解释为什么只采购更多 GPU 或只优化存储都无法解决端到端吞吐。
进入 LLM 时代后，系统重新面临 ``数据移动'' 与 ``数值计算'' 的双重压力。讲者在 37 分钟附近明确提出，基础设施演化可抽象为两条主轴：IO（数据搬运）与 Compute（数值计算）。两者耦合方式决定架构上限。

\begin{importantbox}{AI Cloud 的工程结论}
不是 ``算力越大越好''，而是 ``算力、通信、存储、调度是否协同''。任何单点拉满都会在链路其他位置形成瓶颈。
\end{importantbox}

\begin{warningbox}{盲目复用旧云抽象的风险}
把 AI 训练/推理任务完全当作传统微服务调度，常导致跨机通信抖动、数据加载拥塞和任务碎片化，最终表现为算力利用率低与成本失控。
\end{warningbox}

\subsection{本章小结}
AI Cloud 不是传统云的线性升级，而是回到 ``计算-通信-数据'' 一体化优化问题。系统设计必须从工作负载本身出发。

